\section{Conclusion}\label{sec:conclusion}
This thesis asked whether a post-quantum anonymous credential can run inside a zero-knowledge virtual machine (zkVM) on consumer hardware without breaking the security properties it must provide, with BDEC instantiated with the PLUM signature as the concrete workload on a personal computer (Apple M5~Pro, $24$~GB). The answer has several parts. Precompile-free PLUM verification does not run, because the field-mismatch tax of its large-prime-field algebraic hashing leaves it non-terminating on the target machine (over five hours without terminating; Section~\ref{sec:evaluation}). A custom precompile suite makes it run, turning the non-terminating emulation into a finite proof whose cost can be attributed and bounded. That clears the feasibility bar we can state operationally: completion within the target machine's $24$~GB envelope in bounded wall-clock. Sub-second interactive-showing latency, orders of magnitude below the tens of minutes to hours these proofs take, stays out of reach and is left to future work. With the credential running, a conditional preservation theorem bounds the shift to its unforgeability, anonymity, and unlinkability advantages, and finds them preserved under five explicitly named assumptions and, for the knowledge soundness of the multi-shard receipt, one named extractor conjecture (Section~\ref{sec:security}). One property the solution still leaves out of reach is deployable post-quantum anonymity, because the only zero-knowledge wrap the toolchain ships is pairing-based. The cost criterion, the hash-cost inversion, and the deployment decision rule fill in that answer, and the following subsections take them in turn. We built the construction on SP1 and measured it against an Aurora static-circuit baseline.

\subsection{Main Findings}\label{subsec:conclusion-findings}
\paragraph{The cost criterion.}
A precompile lowers core-proving cost only on the field-mismatch axis, and whether it pays is a trace-area inequality read against two baselines. Against a precompile's own emulated baseline the criterion measures \emph{feasibility}. Against a security-matched control it measures whether the custom chip is worth building. The three shipped modules fall out differently under the two. The Griffin-$\mathbb{F}_p^{192}$ hash pays for feasibility, recovering a finite proof from an emulation stopped unfinished after five hours, yet loses the security-matched comparison, proving slower than a software SHA-3 at matched security ($14.24$ versus $13.24$~min, $n=5$) in an inversion whose sign the trace-area model postdicts. The field-multiplication chip is predicted not to pay, a design-level trace-area comparison rather than an end-to-end measurement, because a dedicated chip's whole-table overhead is not amortised when the host precompile already serves its arithmetic (a dedicated chip is in fact narrower per operation; \S\ref{sec:system}). And the power-residue symbol, irreducibly large-prime yet already-served, still pays $10.33\times$ on total committed trace area. That saving comes from the ${\approx}261{:}1$ control-overhead collapse rather than from any arithmetic saving, and that overhead exists only because the primitive is multi-limb, so this saving too is a field-mismatch effect (Sections~\ref{sec:theoretical} and~\ref{sec:evaluation}).

\paragraph{The dual obstruction.}
The sharpest \emph{finding}, on the security axis rather than the cost axis, is the anonymity--deployability tension (Section~\ref{sec:security}): the dual obstruction to post-quantum anonymity reduces to its single remaining side, the second obstruction. Feasibility is not the issue. The full zero-knowledge wrap chain completes within the memory budget, on consumer hardware and without the default path's multi-day regeneration, so neither the multi-day rebuild the default path would need nor the memory budget is what blocks anonymity (the figures and the reduced verifying-key table are in Section~\ref{sec:evaluation}). What blocks it is a fact about the primitive. The wrap the toolchain ships is pairing-based, so the anonymity it delivers is classical rather than post-quantum, and no proof on this stack is at once zero-knowledge and post-quantum. Recovering post-quantum anonymity needs a post-quantum-sound recursive zero-knowledge wrap, a {STARK}- or lattice-based argument the deployed toolchain does not yet provide. This obstruction holds independently of any timing number and of the open AIR-soundness Assumption~\ref{asm:air} (the dependency direction is recorded in Section~\ref{sec:security} and restated under Limitations below). Pushed to its information-theoretic limit, the boundary is proved: for the hash-committed showing class this construction uses, \emph{everlasting} anonymity cannot be established in the standard model (Theorem~\ref{thm:no-everlasting}), and the conditional route out is a non-interactive statistically-hiding trace commitment (Proposition~\ref{prop:everlasting}).

\paragraph{Feasibility, and two falsified predictions.}
Around that finding, the measurements establish the surrounding cost picture. The custom Griffin-$\mathbb{F}_p^{192}$ precompile turns a non-terminating emulation into a finite measurement: without it, standalone PLUM verification does not complete on the target hardware within the hours attempted (over five hours without terminating; Section~\ref{sec:evaluation}); with it, the verification proves in $14.24$ minutes (mean of five fixed-config runs, $14.16$--$14.44$; Section~\ref{sec:evaluation}) on SP1 at $\lambda=80$, under the documented reduced-memory configuration, default settings exhausting memory (Section~\ref{sec:evaluation}). Two of the measurements falsify natural predictions. First, the algebraic-hash precompile is \emph{slower} than a software SHA-3 control at matched security in every run we took, on both builds we tested. The sign is robust; under fixed-config replication ($n=5$ per arm) the magnitude is ${\approx}1.08\times$ in prove wall-clock ($14.24$ versus $13.24$ minutes). Arithmetising a $199$-bit-field algebraic hash through a custom chip costs more than a hash the substrate already supports. The trace-area accounting of Section~\ref{sec:theoretical} postdicts the sign of this inversion from the chips' committed dimensions: the custom chip adds roughly $5.6\times10^{7}$ committed trace cells per verification (a count derived from the chip's committed dimensions, not a logged trace readout) that the control never pays, a term the guest-cycle accounting alone does not carry. Second, against the transparent, non-preprocessing Aurora baseline the zkVM's update-churn advantage never shows up in wall-clock. The static-side recompile is sub-second ($R_{\mathsf{static}}\approx 0.3$~s, Section~\ref{ssec:churn-eval}), so what separates the two regimes is the regeneration-and-re-audit \emph{footprint} (Table~\ref{tab:churn}), a cost in engineering and assurance rather than in proving time. A wall-clock flexibility advantage would appear only against a \emph{preprocessing} baseline, which is not what BDEC uses.

\paragraph{The credential relations.}
The credential relations themselves run correctly inside the zkVM. Credential generation executes both of its signature verifications to acceptance in $9.12\times10^{10}$ guest cycles, and the presentation relation in $1.36\times10^{11}$ (one credential shown) and $1.82\times10^{11}$ (two) cycles, consistent with the additive cost model. On SP1 both relations complete a succinct, non-zero-knowledge prove (credential generation in $68.96$~min ($1.47$~GB) under the config-verified reproducible configuration, presentation in $103.34$~min ($2.2$~GB) at $k{=}1$ and $136.82$~min ($2.9$~GB) at $k{=}2$; Section~\ref{sec:evaluation}). The zero-knowledge wrap of the credential relation is produced on consumer hardware as well: credential generation wraps end-to-end in $4.20$~hours within the memory budget, peaking near $15$~GiB, so the credential-generation proof, like the standalone verification, completes. The classical anonymity it then provides is conditional on the open key-privacy assumption (Assumption~\ref{asm:keypriv}). Both relations moreover bind their public statement to the receipt, demonstrated end to end, and for issuance that binding and the zero-knowledge wrap combine in one measured proof (a bound CreGen wrap in $4.18$~h; Section~\ref{sec:evaluation}), so the deployable object is a statement-bound zero-knowledge credential rather than two separate demonstrations. The presentation wrap completes and verifies at both $k=1$ ($11.76$~hours end-to-end, plain wrap) and $k=2$ ($10.21$~hours, statement-bound; the two are different configurations, not a matched $k$-series), each within the memory budget and peaking near $15$~GiB ($14.93$ and $15.66$~GiB). The core-proof shard count ($1516\to2011$) and peak memory grow with the credential count while staying within the $24$~GB budget. What blocks post-quantum anonymity is separate, and independent of $k$. The only zero-knowledge wrap the toolchain ships is pairing-based (the second obstruction).

\subsection{Answer to the Research Question}\label{subsec:conclusion-answer}
The narrower, mechanism-level question this thesis poses is \emph{when a precompile pays}, and the answer is a trace-area inequality read against two baselines. Against a precompile's own emulated baseline it pays exactly when the core trace area it removes exceeds the chip area it adds; by that feasibility test the Griffin chip pays, recovering a finite proof from an emulation stopped unfinished after five hours. Against the best available alternative it need not, and here it does not: the Griffin precompile proves slower than a software SHA-3 control that commits no chip area at all. The research question asks the larger thing, whether a post-quantum anonymous credential runs inside a zkVM without breaking its security, and its answer is \emph{not yet}, for a reason more specific than ``too slow.'' Standalone PLUM verification becomes feasible on consumer hardware, at $\lambda=80$ on SP1 under tuned memory settings and conditional on Assumption~\ref{asm:air}, once the precompile absorbs the field-mismatch tax. Deployable post-quantum \emph{anonymity} is not reached, because it requires a zero-knowledge wrap that is at once feasible to generate within $24$~GB and post-quantum-sound. The wrap the pipeline exposes is now feasible, running to completion within budget, but it is pairing-based and therefore not post-quantum-sound, for the reason of the previous subsection. The substrate choice is therefore governed by the deployment-lifecycle decision rule of Section~\ref{sec:discussion}: where the verification predicate is fixed and proofs are frequent, the specialised static circuit is preferable; where the predicate, signature, or security parameters evolve over the deployment's lifetime, the zkVM's fixed universal relation confines each change to a guest-program edit, provided the change leaves the precompiled primitive suite intact, since a hash-family or field-width change regenerates the Griffin AIR. Against a non-preprocessing baseline its advantage there is a smaller re-audit footprint, not a faster runtime.

\subsection{Limitations}\label{subsec:conclusion-limitations}
The results are bounded in ways we state explicitly.

\paragraph{Parameter level.}
Prove-mode wall-clock is measured at $\lambda=80$, the largest level at which it is observable within the memory budget and not a deployment security level; the field-multiplication ablation is reported at $\lambda=128$, while the Griffin-precompile cycle ratio ($58.5\times$) and the credential relation cycle counts are execute-mode figures at $\lambda=80$.

\paragraph{The static baseline.}
The static baseline is Loquat over $\mathbb{F}_{p_{127}}$ rather than same-scheme PLUM, with the per-proof figure measured at zero knowledge disabled and the zero-knowledge-enabled counterpart proving in about $22$ minutes (mean of three runs, $21.4$--$22.9$) (Section~\ref{sec:evaluation}), so the per-proof comparison is indicative and the lifecycle crossover is bounded in form but not located in wall-clock.

\paragraph{Run counts.}
The standalone PLUM Cell~2 and Cell~3 prove figures are means of five fixed-config runs (ranges reported); the credential-relation prove figures are single runs (Section~\ref{sec:evaluation}).

\paragraph{Conditional security.}
The security guarantees are conditional: the constraint-level soundness of the Griffin AIR (Assumption~\ref{asm:air}), the existence of a PLUM signature-transcript simulator (Assumption~\ref{asm:szk}), the cross-table lookup-argument binding (Assumption~\ref{asm:lookup}), PLUM signature key-privacy for issuance anonymity (Assumption~\ref{asm:keypriv}), and the quantum-random-oracle lift (Assumption~\ref{asm:qrom}) are carried as the five essential open assumptions (the one further named premise of Section~\ref{sec:security}, the black-box-transfer of Assumption~\ref{asm:bbt}, is inherited or standard), together with the multi-shard extractor conjecture (Conjecture~\ref{conj:extract-recursion}) on which the composite receipt's knowledge soundness rests, and with the canonical prime instantiation a separate premise (Remark~\ref{rem:prime-substitution}). Of these, Assumption~\ref{asm:air} conditions only the positive results, the $14.24$-minute feasibility figure and the security-preservation theorem; the negative findings, including the surviving obstruction (the second), the SHA-3 inversion, the precompile-free baseline's non-completion within the time budget, and the absence of a wall-clock flexibility advantage, would survive or strengthen under an unsound AIR (the argument, that an unsound AIR only understates the Griffin-arm cost, is in Section~\ref{sec:security}).

\paragraph{Statement binding and wrap run-counts.}
The wrap figures are single runs; that is the limitation this paragraph carries, and the positive record it bounds is as follows. The SP1 prototype binds the receipt journal to the public statement for both credential relations, demonstrated end-to-end by two cross-issuer issuances and a $k=2$ presentation, each accepting with its committed statement recovered (Section~\ref{sec:evaluation}); a single statement-bound and zero-knowledge-wrapped run is measured for credential generation, a bound CreGen wrap completing and verifying in $4.18$~h at $15.2$~GiB peak, essentially the plain wrap's cost, and the $k=2$ showing counterpart likewise, a statement-bound wrap completing and verifying in $10.21$~h at a $15.66$~GiB peak (Section~\ref{sec:evaluation}). The wrap's $24$~GB-sufficiency is established by completed end-to-end runs of the PLUM-verification, credential-generation, and both the $k=1$ and $k=2$ presentation wraps, each peaking near $15$~GiB (Section~\ref{sec:evaluation}).
\subsection{Future Work}\label{subsec:conclusion-future}

\paragraph{The next measurement.}
The most informative next measurement is a \emph{same-scheme} per-proof comparison, running one scheme on both substrates at matched security: a $199$-bit $\mathbb{F}_p$ Aurora harness to run PLUM statically (the current harness is $127$-bit only), fixing the sign of the per-proof crossover that the decision rule of Section~\ref{sec:discussion} currently leaves as a conjecture.

\paragraph{Security of the SHA-3 substitution.}
The SHA-3 arm of Section~\ref{sec:evaluation} is a diagnostic control, and its security is deliberately unexamined: PLUM's analysis covers its own algebraic-hash instantiation, and a bit-oriented hash is ruled out by cost in PLUM's native circuit setting, where hash evaluations already account for ${\approx}91\%$ of the verification constraints (Section~\ref{sec:theoretical}), before any security question arises. On the zkVM substrate that cost objection disappears, so the question becomes live. Settling it would require re-specifying the encoding of field elements into a bit-oriented sponge, the challenge-expansion sampling, and the Merkle commitment instantiation, and re-deriving the soundness bounds of the underlying IOP under those choices. If the substituted scheme proves secure, the inversion of Section~\ref{sec:evaluation} becomes deployment-relevant rather than diagnostic, and the case for algebraic hashes on this substrate narrows further. We leave the analysis as future work.

\paragraph{The open assumptions.}
The open assumptions are the substantive cryptographic work that remains: Assumption~\ref{asm:air}, whose per-row layer carries a written proof (Appendix~\ref{app:proofs}) and rejection-direction fault-injection coverage for seven of the ten constraint families, by extending the suite through a full-prover harness to the lookup-borne families (round count, memory binding, cross-row threading) and by machine-checking the code-to-constraint transcription, for instance by the method of Arguzz~\cite{hochrainer2025arguzz}; Assumption~\ref{asm:szk} by constructing a PLUM signature-transcript simulator; and the post-quantum-anonymity goal by a post-quantum-sound zero-knowledge wrap (a STARK- or lattice-based recursive argument) in place of the pairing-based wrap the toolchain currently exposes.

\paragraph{A post-quantum-anonymity path.}
The post-quantum-anonymity obstruction of Section~\ref{ssec:zk-gap} is largely an engineering gap plus one open lemma. The classical zero-knowledge the produced wrap already supplies can be made post-quantum without a pairing-based wrap, by masking the prover's own low-degree test rather than wrapping its output: for a {FRI}/{STARK}-style prover by masked {FRI}~\cite{cryptoeprint:2024/1037}, and for the multilinear Basefold-style commitment of the SP1 prover by the VEIL compiler~\cite{cryptoeprint:2026/683}, both plausibly post-quantum and neither reintroducing a quantum-vulnerable assumption. VEIL exists as code (it ships in the SP1 fork, as yet unwired), but wiring it is not a flag-flip. Its one-mask-per-commitment discipline collides with the repeated openings of the multi-shard SP1 prover, and the honest-verifier zero-knowledge of the masked low-degree test at the deployed instantiation, a per-round leaf-masking statement for STIR/WHIR at $\eta=4$ over the $199$-bit field, is an open lemma we scope but do not prove here, the per-round-committed-oracle analog of the single-quotient statement of masked {FRI}~\cite{cryptoeprint:2024/1037}. What remains is therefore the reconciliation of VEIL's opening discipline with multi-shard proving, and then measuring whether the resulting post-quantum-sound prover stays within the $24$~GB budget the pairing-based wrap already meets. Closing the obstruction this way would deliver \emph{computational} post-quantum anonymity in the quantum random-oracle model. It would not deliver \emph{everlasting} anonymity: for the hash-committed showings, everlasting anonymity cannot be established in the standard model (Theorem~\ref{thm:no-everlasting}), and reaching it would require a statistically-hiding trace commitment, a separate direction we leave open.

\paragraph{A precompile generator.}
A further direction concerns the precompile construction itself. The three chips of Section~\ref{sec:system} instantiate one interface by hand, and the cost criterion decides, per primitive, which of them is worth building. The natural generalisation is a generator that emits a sound precompile directly from a primitive's specification: given the S-box degree, state width, round structure, and field of a large-prime algebraic primitive, it would produce the AIR together with a machine-checkable soundness obligation from a single meta-theorem, rather than writing and auditing each chip separately. This would turn the interface of Section~\ref{ssec:soundness} from a contract instantiated three times into a reusable construction method, and it would let the criterion drive an automated build-or-skip decision in the manner of the autoprecompile line~\cite{powdr2025autoprecompiles}, but over native large-field arithmetic rather than merged emulation. The soundness meta-theorem is the hard part, and we leave it as future work.

\subsection{Closing}\label{subsec:conclusion-closing}
Post-quantum anonymous credential systems require both cryptographic soundness and operational evolvability, and this thesis measured the price of pursuing both on the same substrate. Its contribution is to map that frontier precisely. The precompile makes PLUM signature verification at the $80$-bit level finite to prove on a personal computer. The deployment-lifecycle separation between the static and zkVM regimes is real, though it surfaces as a re-audit footprint and stays out of the wall-clock numbers. The dual obstruction to deployable post-quantum anonymity reduces to its single remaining obstruction, the second: the anonymity-providing wrap completes within the memory budget, but it is pairing-based and therefore classical rather than post-quantum, and no amount of additional memory, favourable timing, or resolution of Assumption~\ref{asm:air} makes it post-quantum. That obstruction, named precisely and with its cost accounted for, is the main result this thesis leaves for the design of the next such system.
